\section{Reproducibility}
\label{app:protocol}
\subsection{Graph preprocessing and fixed checkpoints}
Table~\ref{tab:datasets} records the main dataset dimensions. Existing graph tensors and preprocessing are retained, but official masks are no longer used. Each task follows the original ProG implementation with two independent permutations for the 90\% test set and 10\% support pool. The first $k$ members of each class are selected; other nodes provide unlabeled context. Support/test overlap is retained, with indices and counts stored in the result package.
\begin{table}[H]
\centering
\linespread{1}\selectfont
\captionsetup{font=small,skip=5pt}
\caption{\lang{\textbf{Dataset and prompt sizes.}}{\textbf{数据集规模与提示参数量。}}}
\label{tab:datasets}
\fontsize{9}{11}\selectfont
\setlength{\tabcolsep}{1pt}
\renewcommand{\arraystretch}{1.18}
\setlength{\heavyrulewidth}{0.8pt}
\setlength{\lightrulewidth}{0.4pt}
\begin{tabular*}{\linewidth}{@{\extracolsep{\fill}}lrrrr@{}}
\toprule
\lang{Dataset}{数据集} & \lang{Nodes}{节点} & \lang{Features}{特征} & \lang{Classes}{类别} & \lang{\shortstack{Prompt\\params.}}{\shortstack{提示\\参数}} \\
\midrule
Cora & 2,708 & 1,433 & 7 & 896 \\
CiteSeer & 3,327 & 3,703 & 6 & 768 \\
PubMed & 19,717 & 500 & 3 & 384 \\
\bottomrule
\end{tabular*}
\end{table}


The encoder consists of two GCNConv layers of width 128 with LeakyReLU between them. The first layer includes the pretrained affine transformation and graph aggregation; the prompt is added to its output before the nonlinearity. The second convolution has no additional activation in the evaluated encoder. The same checkpoint is used by every comparison method on a dataset. Checkpoint identity is checked at preparation, and task indices are retained. The recorded pretraining family is \texttt{Edgepred\_Gprompt}; the pretraining random seed is unavailable.

\subsection{Training and prediction details}
The main text describes the active forward computation. For reproducibility, Adam uses $\beta_1=0.9$, $\beta_2=0.999$, $\epsilon=10^{-8}$, learning rate 0.01, and weight decay zero. The prompt regularizer is the mean of squared token entries multiplied by $10^{-3}$. It is separate from optimizer weight decay. The channel scale uses the sample standard deviation, with denominator $n-1$. Euclidean normalization uses the implementation's $10^{-12}$ norm floor. There is no distillation or additional test-label objective; overlapping test nodes receive supervision as support nodes.

For the support loss, let $\bar z_i=\norm{Z_i}$ and $m_c=\sum_{j\in S_c}\bar z_j$. When $\min_c k_c>1$, the normalized prototype used for training support node $i$ is
\begin{equation}
\begin{aligned}
p_c^{(-i)}&=\norm{\frac{m_c-\mathbb I[y_i=c]\bar z_i}{k_c-\mathbb I[y_i=c]}},\\
\ell_{ic}&=\bar z_i^\top p_c^{(-i)}/0.2.
\end{aligned}
\label{eq:loo}
\end{equation}
When $\min_c k_c=1$, the inclusive class prototype $\norm{m_c/k_c}$ replaces $p_c^{(-i)}$. During prediction, define
\begin{equation}
\begin{aligned}
z_i^*&=\norm{Z_i-\frac1n\sum_{j\in V}Z_j},\\
p_c^*&=\norm{\frac1{k_c}\sum_{i\in S_c}z_i^*},\\
\widehat y_i&=\arg\max_c (z_i^*)^\top p_c^*.
\end{aligned}
\label{eq:predict}
\end{equation}
The probability scores for all five metrics are obtained by applying softmax to cosine logits divided by 0.2. Training does not apply the output centering used at prediction. The readout is a fixed implementation choice, not an experimental factor in this paper.

\paragraph{Procedure.}
Given a graph, a fixed checkpoint, and support labels:
\begin{enumerate}
\item Normalize adjacency and compute Eq.~\eqref{eq:attributes} and the first-layer preactivation.
\item Initialize support prototypes with Eq.~\eqref{eq:anchor} and apply Eqs.~\eqref{eq:assignment}--\eqref{eq:refinement} three times.
\item Compute the fixed route with Eq.~\eqref{eq:routing}; for the shared-prompt ablation, replace it by $B_{ic}=1/C$.
\item Initialize $P=0$ and cache the hidden-channel scales. Keep the encoder and route fixed.
\item For 100 steps, run the prompted encoder, form support logits, and update $P$ using Eq.~\eqref{eq:loss}.
\item Compute final prototype scores using Eq.~\eqref{eq:predict}; seal predictions before scoring with query labels.
\end{enumerate}

\subsection{Baseline fitting and seed control}
Task seeds are 30093000 through 30093019 and training seeds are 0,1,2. Each task first draws test nodes, then independently draws the support pool, selects the first $k$ nodes per class, and shuffles support order. Methods share task identities and shot-specific support sets are nested. Insufficient classes retain $k_c<k$ without replenishment. The no-training Frozen method retains three seed slots for alignment; these are not independent fits.

Local Full FT uses support cross-entropy and a task readout. Its parameters when the minimum class support count is one are learning rate 0.01 and weight decay 0.0001. When the minimum count exceeds one, each of $\min(3,\min_c k_c)$ support-only folds holds out one node per class. Four candidate configurations are ranked by mean held-out support cross-entropy, then Full FT is refit on the complete support set. No external validation-label budget is used. This local protocol controls labeled information and final step count, but not total selection compute or readout parameterization. It does not describe the imported literature baselines.

\subsection{Condition-specific statistics}
For a fixed dataset, shot count, checkpoint, and method, let $a_{st}$ denote a metric with task seed $s$ and training seed $t$. We compute
\begin{equation}
\begin{aligned}
\bar a_s&=\frac13\sum_{t=1}^{3}a_{st},\qquad
\widehat a=\frac1{20}\sum_{s=1}^{20}\bar a_s,\\
\widehat\sigma&=\sqrt{\frac1{19}\sum_{s=1}^{20}(\bar a_s-\widehat a)^2}.
\end{aligned}
\label{eq:stats}
\end{equation}
Local cells report $100\widehat a\pm100\widehat\sigma$, not the standard error. For local paired comparisons, $d_s=\bar a_s^{\mathrm{PRPrompt}}-\bar a_s^{\mathrm{reference}}$ is resampled as a task-draw block. The percentile bounds of 10,000 resampled means form the reported 95\% interval. All intervals are recomputed for this rerun. Literature cells retain their source five-trial mean and standard deviation; Eq.~\eqref{eq:stats} does not apply to them. No metric average or mean effect is computed across datasets, shots, methods, or sources.

Accuracy, Macro-F1, Micro-F1, macro one-vs-rest ROC-AUC, and macro average precision are retained in the data package. In this single-label multiclass setting, Micro-F1 coincides with Accuracy; it is included for completeness. Support and test sets are redrawn across tasks; intervals reflect task sampling and training variation conditional on the graph, not independent graph samples.

\subsection{Environment, provenance, and numerical qualification}
The existing ProG environment is reused with PyTorch 2.11.0+cu130 and the same graph inputs and checkpoints. Per-unit fitting times are recorded; concurrent wall-clock times do not establish a method-level efficiency advantage.

The eight-dataset rerun completed 5760 formal units and 96 separate two-step smoke units. Local checks cover coverage, finite metrics, paired support/test indices, and recomputed means, standard deviations, and bootstrap intervals. Saved-embedding readout uses a $2\times10^{-6}$ tolerance. An old strict filter row-sum assertion was converted to error recording without changing operator arithmetic. No exhaustive file-hash scan or fresh CUDA bitwise audit was repeated; earlier numerical audits are not presented as evidence for this rerun.

External Accuracy values are transcribed from \citet{mtg2026}, Tables 4--6 (PDF pages 37--38 of arXiv:2606.03290v1). The record preserves table/page identifiers, decimal strings, source URL and PDF hash. It contains 147 node-classification cells excluding MTG; the manuscript uses 63 cells for its three datasets. Source tables combine ProG results and author reproductions, without cell-specific attribution or selected pretraining identifiers. Raw predictions are unavailable in the imported record, so no external paired tests are computed. Local historical prompt measurements remain archived but are not substituted for missing literature metrics.

Code and results are available at \url{https://github.com/Mes-sala/PRPrompt}; execution scripts are in \texttt{experiments/prog\_split\_20261002}. This work reruns adaptation and statistics under original ProG sampling without new pretraining. The first seven datasets reuse their original Edgepred\_Gprompt checkpoints with unknown pretraining seeds. ogbn-arxiv reuses its previously trained full-graph edge-prediction checkpoint with pretraining seed 42. Weights, graph snapshots, and model states remain on the server; local artifacts contain code, predictions, splits, tables, and documents.

\clearpage
\twocolumn[{
\section{Additional Metrics}
\label{app:more}
The following tables extend the main comparison with the remaining stored metrics under the same seed aggregation rule. Every column is a separate dataset/shot condition.
\par\medskip\noindent\begin{minipage}{\textwidth}
\centering
\linespread{1}\selectfont
\captionsetup{font=small,skip=5pt}
\captionof{table}{\lang{\textbf{Micro-F1 on few-shot node classification.} Higher is better ($\uparrow$).}{\textbf{少样本节点分类的 Micro-F1。}越高越好（$\uparrow$）。}}
\label{tab:main-micro_f1}
\fontsize{9}{11}\selectfont
\setlength{\tabcolsep}{1pt}
\renewcommand{\arraystretch}{1.18}
\setlength{\heavyrulewidth}{0.8pt}
\setlength{\lightrulewidth}{0.4pt}
\begin{tabular*}{\linewidth}{@{\extracolsep{\fill}}l*{9}{r}@{}}
\toprule
\lang{Method}{方法} & \multicolumn{3}{c}{\textbf{Cora}} & \multicolumn{3}{c}{\textbf{CiteSeer}} & \multicolumn{3}{c}{\textbf{PubMed}} \\
\cmidrule(lr){2-4}\cmidrule(lr){5-7}\cmidrule(l){8-10}
 & 1-shot & 3-shot & 5-shot & 1-shot & 3-shot & 5-shot & 1-shot & 3-shot & 5-shot \\
\midrule
\frozenname & \meanstd{24.64}{4.63} & \meanstd{30.86}{3.27} & \meanstd{35.16}{2.40} & \meanstd{21.69}{3.47} & \meanstd{24.96}{3.51} & \meanstd{27.52}{3.72} & \secondmeanstd{53.72}{10.60} & \secondmeanstd{62.49}{9.18} & \secondmeanstd{65.68}{5.56} \\
\ftname & \secondmeanstd{45.15}{8.56} & \secondmeanstd{59.29}{6.26} & \secondmeanstd{66.12}{5.09} & \secondmeanstd{33.59}{5.82} & \secondmeanstd{42.05}{6.91} & \secondmeanstd{49.92}{5.85} & \meanstd{49.90}{8.69} & \meanstd{58.72}{8.45} & \meanstd{64.14}{7.08} \\
\midrule\rowcolor{black!5}
\textbf{\method} & \bestmeanstd{53.80}{8.45} & \bestmeanstd{69.97}{3.53} & \bestmeanstd{72.87}{3.36} & \bestmeanstd{43.92}{11.01} & \bestmeanstd{59.67}{5.98} & \bestmeanstd{64.70}{3.53} & \bestmeanstd{59.77}{7.47} & \bestmeanstd{65.09}{8.21} & \bestmeanstd{68.34}{6.97} \\
\bottomrule
\end{tabular*}
\par\vspace{3pt}
\begin{minipage}{\linewidth}\fontsize{8}{9.5}\selectfont\raggedright \lang{Local original-ProG runs with support/test overlap: mean $\pm$ sample std. over 20 tasks after averaging three training seeds per task. \textbf{Bold}/\underline{underline}: best/second in each condition. No cross-condition average.}{本地原ProG实现，保留支持／测试重叠：先平均每任务3个训练seed，再报告20个任务的均值 $\pm$ 样本标准差。\textbf{粗体}／\underline{下划线}表示同条件最优／次优；不跨条件平均。}\lang{ Literature baselines do not report this metric in Tables 4--6.}{文献表 4--6 未报告这些基线的本指标。}\end{minipage}
\end{minipage}\par

\par\medskip\noindent\begin{minipage}{\textwidth}
\centering
\linespread{1}\selectfont
\captionsetup{font=small,skip=5pt}
\captionof{table}{\lang{\textbf{Macro ROC-AUC on few-shot node classification.} Higher is better ($\uparrow$).}{\textbf{少样本节点分类的 Macro ROC-AUC。}越高越好（$\uparrow$）。}}
\label{tab:main-roc_auc_macro}
\fontsize{9}{11}\selectfont
\setlength{\tabcolsep}{1pt}
\renewcommand{\arraystretch}{1.18}
\setlength{\heavyrulewidth}{0.8pt}
\setlength{\lightrulewidth}{0.4pt}
\begin{tabular*}{\linewidth}{@{\extracolsep{\fill}}l*{9}{r}@{}}
\toprule
\lang{Method}{方法} & \multicolumn{3}{c}{\textbf{Cora}} & \multicolumn{3}{c}{\textbf{CiteSeer}} & \multicolumn{3}{c}{\textbf{PubMed}} \\
\cmidrule(lr){2-4}\cmidrule(lr){5-7}\cmidrule(l){8-10}
 & 1-shot & 3-shot & 5-shot & 1-shot & 3-shot & 5-shot & 1-shot & 3-shot & 5-shot \\
\midrule
\frozenname & \meanstd{63.21}{4.73} & \meanstd{69.93}{2.36} & \meanstd{73.21}{2.16} & \meanstd{56.08}{3.35} & \meanstd{59.27}{3.49} & \meanstd{61.37}{2.97} & \meanstd{71.26}{10.76} & \secondmeanstd{80.38}{8.27} & \meanstd{83.07}{4.42} \\
\ftname & \secondmeanstd{79.66}{4.17} & \secondmeanstd{88.47}{3.03} & \secondmeanstd{91.72}{2.18} & \secondmeanstd{69.34}{5.25} & \secondmeanstd{74.01}{3.87} & \secondmeanstd{78.59}{3.00} & \secondmeanstd{71.68}{7.97} & \meanstd{79.93}{5.25} & \bestmeanstd{83.87}{4.51} \\
\midrule\rowcolor{black!5}
\textbf{\method} & \bestmeanstd{81.14}{5.37} & \bestmeanstd{90.82}{2.12} & \bestmeanstd{92.22}{1.34} & \bestmeanstd{72.42}{7.93} & \bestmeanstd{82.11}{2.86} & \bestmeanstd{84.12}{2.17} & \bestmeanstd{77.21}{6.53} & \bestmeanstd{81.72}{7.67} & \secondmeanstd{83.84}{5.88} \\
\bottomrule
\end{tabular*}
\par\vspace{3pt}
\begin{minipage}{\linewidth}\fontsize{8}{9.5}\selectfont\raggedright \lang{Local original-ProG runs with support/test overlap: mean $\pm$ sample std. over 20 tasks after averaging three training seeds per task. \textbf{Bold}/\underline{underline}: best/second in each condition. No cross-condition average.}{本地原ProG实现，保留支持／测试重叠：先平均每任务3个训练seed，再报告20个任务的均值 $\pm$ 样本标准差。\textbf{粗体}／\underline{下划线}表示同条件最优／次优；不跨条件平均。}\lang{ Literature baselines do not report this metric in Tables 4--6.}{文献表 4--6 未报告这些基线的本指标。}\end{minipage}
\end{minipage}\par

\par\medskip\noindent\begin{minipage}{\textwidth}
\centering
\linespread{1}\selectfont
\captionsetup{font=small,skip=5pt}
\captionof{table}{\lang{\textbf{Macro AP on few-shot node classification.} Higher is better ($\uparrow$).}{\textbf{少样本节点分类的 Macro AP。}越高越好（$\uparrow$）。}}
\label{tab:main-average_precision_macro}
\fontsize{9}{11}\selectfont
\setlength{\tabcolsep}{1pt}
\renewcommand{\arraystretch}{1.18}
\setlength{\heavyrulewidth}{0.8pt}
\setlength{\lightrulewidth}{0.4pt}
\begin{tabular*}{\linewidth}{@{\extracolsep{\fill}}l*{9}{r}@{}}
\toprule
\lang{Method}{方法} & \multicolumn{3}{c}{\textbf{Cora}} & \multicolumn{3}{c}{\textbf{CiteSeer}} & \multicolumn{3}{c}{\textbf{PubMed}} \\
\cmidrule(lr){2-4}\cmidrule(lr){5-7}\cmidrule(l){8-10}
 & 1-shot & 3-shot & 5-shot & 1-shot & 3-shot & 5-shot & 1-shot & 3-shot & 5-shot \\
\midrule
\frozenname & \meanstd{24.76}{4.19} & \meanstd{30.07}{2.96} & \meanstd{34.33}{2.50} & \meanstd{21.33}{2.42} & \meanstd{23.68}{2.88} & \meanstd{25.30}{3.00} & \meanstd{55.51}{10.07} & \secondmeanstd{66.73}{9.18} & \meanstd{68.93}{6.61} \\
\ftname & \secondmeanstd{50.83}{7.37} & \secondmeanstd{67.63}{5.63} & \secondmeanstd{75.34}{4.11} & \secondmeanstd{36.52}{6.95} & \secondmeanstd{43.52}{6.52} & \secondmeanstd{50.83}{5.30} & \secondmeanstd{56.10}{9.47} & \meanstd{66.58}{7.45} & \bestmeanstd{72.23}{6.16} \\
\midrule\rowcolor{black!5}
\textbf{\method} & \bestmeanstd{56.27}{8.10} & \bestmeanstd{74.43}{4.49} & \bestmeanstd{78.20}{2.98} & \bestmeanstd{42.17}{10.92} & \bestmeanstd{57.10}{5.18} & \bestmeanstd{61.42}{3.29} & \bestmeanstd{61.55}{8.90} & \bestmeanstd{68.29}{9.31} & \secondmeanstd{70.78}{8.65} \\
\bottomrule
\end{tabular*}
\par\vspace{3pt}
\begin{minipage}{\linewidth}\fontsize{8}{9.5}\selectfont\raggedright \lang{Local original-ProG runs with support/test overlap: mean $\pm$ sample std. over 20 tasks after averaging three training seeds per task. \textbf{Bold}/\underline{underline}: best/second in each condition. No cross-condition average.}{本地原ProG实现，保留支持／测试重叠：先平均每任务3个训练seed，再报告20个任务的均值 $\pm$ 样本标准差。\textbf{粗体}／\underline{下划线}表示同条件最优／次优；不跨条件平均。}\lang{ Literature baselines do not report this metric in Tables 4--6.}{文献表 4--6 未报告这些基线的本指标。}\end{minipage}
\end{minipage}\par

}]
\clearpage
\twocolumn[{
\section{Paired Differences and Uncertainty}
\label{app:paired}
Each cell reports \method{} minus a local reference in percentage points, with its unadjusted 95\% paired task-bootstrap interval. All entries preserve original paired statistics. Literature baselines are excluded because the runs are not paired. Tables group conditions for readability but do not average them.
\par\medskip
\par\medskip\noindent\begin{minipage}{\textwidth}
\centering
\linespread{1}\selectfont
\captionsetup{font=small,skip=5pt}
\captionof{table}{\lang{\textbf{Paired Accuracy gains of \method.} Differences in percentage points.}{\textbf{\method 的 Accuracy 配对增益。}差值单位为百分点。}}
\label{tab:paired-accuracy}
\fontsize{9}{11}\selectfont
\setlength{\tabcolsep}{1pt}
\renewcommand{\arraystretch}{1.18}
\setlength{\heavyrulewidth}{0.8pt}
\setlength{\lightrulewidth}{0.4pt}
\begin{tabular*}{\linewidth}{@{\extracolsep{\fill}}l*{9}{r}@{}}
\toprule
\lang{Reference}{对照方法} & \multicolumn{3}{c}{\textbf{Cora}} & \multicolumn{3}{c}{\textbf{CiteSeer}} & \multicolumn{3}{c}{\textbf{PubMed}} \\
\cmidrule(lr){2-4}\cmidrule(lr){5-7}\cmidrule(l){8-10}
 & 1-shot & 3-shot & 5-shot & 1-shot & 3-shot & 5-shot & 1-shot & 3-shot & 5-shot \\
\midrule
\renewcommand{\arraystretch}{1.35}
\frozenname & \pairedcell{+31.42}{28.20}{34.59} & \pairedcell{+41.34}{39.63}{43.04} & \pairedcell{+39.46}{37.23}{41.61} & \pairedcell{+25.98}{23.72}{28.21} & \pairedcell{+35.91}{33.05}{38.39} & \pairedcell{+38.78}{37.37}{40.18} & \pairedcell{+3.26}{0.40}{6.22} & \pairedcell{+3.11}{0.38}{6.27} & \pairedcell{+2.61}{1.13}{4.22} \\[2pt]
\ftname & \pairedcell{+9.17}{6.58}{11.87} & \pairedcell{+9.22}{6.75}{11.79} & \pairedcell{+5.52}{3.73}{7.50} & \pairedcell{+11.36}{9.57}{13.18} & \pairedcell{+18.98}{16.78}{20.94} & \pairedcell{+18.28}{16.22}{20.13} & \pairedcell{+9.93}{6.77}{13.14} & \pairedcell{+8.90}{5.89}{12.26} & \pairedcell{+7.02}{4.96}{9.14} \\[2pt]
\midrule
GPF & \pairedcell{+39.38}{35.60}{42.98} & \pairedcell{+46.44}{44.21}{48.74} & \pairedcell{+44.86}{42.52}{47.57} & \pairedcell{+30.02}{27.04}{33.00} & \pairedcell{+40.83}{38.11}{43.06} & \pairedcell{+44.77}{43.32}{46.16} & \pairedcell{+22.40}{18.82}{26.24} & \pairedcell{+17.37}{14.16}{21.02} & \pairedcell{+15.15}{13.30}{17.01} \\[2pt]
GPF-plus & \pairedcell{+23.78}{19.96}{27.47} & \pairedcell{+30.55}{28.29}{32.52} & \pairedcell{+19.88}{17.62}{22.16} & \pairedcell{+25.71}{22.89}{28.45} & \pairedcell{+32.27}{30.45}{34.03} & \pairedcell{+34.35}{33.05}{35.65} & \pairedcell{+17.57}{13.28}{21.86} & \pairedcell{+11.20}{9.47}{13.11} & \pairedcell{+9.61}{7.84}{11.49} \\[2pt]
\gpromptname & \pairedcell{+32.70}{29.19}{36.29} & \pairedcell{+43.01}{41.38}{44.73} & \pairedcell{+41.66}{40.03}{43.20} & \pairedcell{+27.59}{24.97}{30.09} & \pairedcell{+36.37}{33.26}{38.82} & \pairedcell{+38.51}{37.14}{39.90} & \pairedcell{+8.45}{3.60}{13.47} & \pairedcell{+6.19}{4.04}{8.30} & \pairedcell{+4.03}{1.94}{6.36} \\[2pt]
\midrule
\sharedname & \pairedcell{+31.96}{28.59}{35.21} & \pairedcell{+39.06}{36.49}{41.44} & \pairedcell{+36.94}{34.87}{39.06} & \pairedcell{+26.40}{24.07}{28.73} & \pairedcell{+37.18}{34.90}{39.33} & \pairedcell{+40.45}{39.06}{41.84} & \pairedcell{+4.85}{1.26}{8.26} & \pairedcell{+5.69}{3.25}{8.27} & \pairedcell{+3.02}{1.39}{4.51} \\[2pt]
\bottomrule
\end{tabular*}
\par\vspace{3pt}
\begin{minipage}{\linewidth}\fontsize{8}{9.5}\selectfont\raggedright \lang{Each cell: PRPrompt minus the reference (top); unadjusted 95\% paired support-bootstrap interval (below). All values preserve the original condition-wise statistics.}{每格上行为 PRPrompt 减去对照的差值，下行为未经多重比较校正的 95\% 配对支持 bootstrap 区间。数值沿用原始逐条件统计。}\end{minipage}
\end{minipage}\par


\par\medskip\noindent\begin{minipage}{\textwidth}
\centering
\linespread{1}\selectfont
\captionsetup{font=small,skip=5pt}
\captionof{table}{\lang{\textbf{Paired Macro-F1 gains of \method.} Differences in percentage points.}{\textbf{\method 的 Macro-F1 配对增益。}差值单位为百分点。}}
\label{tab:paired-macro_f1}
\fontsize{9}{11}\selectfont
\setlength{\tabcolsep}{1pt}
\renewcommand{\arraystretch}{1.18}
\setlength{\heavyrulewidth}{0.8pt}
\setlength{\lightrulewidth}{0.4pt}
\begin{tabular*}{\linewidth}{@{\extracolsep{\fill}}l*{9}{r}@{}}
\toprule
\lang{Reference}{对照方法} & \multicolumn{3}{c}{\textbf{Cora}} & \multicolumn{3}{c}{\textbf{CiteSeer}} & \multicolumn{3}{c}{\textbf{PubMed}} \\
\cmidrule(lr){2-4}\cmidrule(lr){5-7}\cmidrule(l){8-10}
 & 1-shot & 3-shot & 5-shot & 1-shot & 3-shot & 5-shot & 1-shot & 3-shot & 5-shot \\
\midrule
\renewcommand{\arraystretch}{1.35}
\frozenname & \pairedcell{+30.84}{27.77}{33.84} & \pairedcell{+40.32}{38.64}{42.08} & \pairedcell{+38.97}{36.88}{40.99} & \pairedcell{+23.87}{21.80}{25.77} & \pairedcell{+33.87}{31.34}{36.14} & \pairedcell{+36.58}{35.22}{37.95} & \pairedcell{+3.93}{0.86}{7.00} & \pairedcell{+5.02}{2.27}{8.44} & \pairedcell{+4.12}{2.53}{5.91} \\[2pt]
\ftname & \pairedcell{+8.97}{6.40}{11.77} & \pairedcell{+8.44}{6.12}{10.90} & \pairedcell{+5.05}{3.34}{6.86} & \pairedcell{+13.27}{11.17}{15.31} & \pairedcell{+19.76}{17.45}{22.00} & \pairedcell{+18.42}{16.07}{20.60} & \pairedcell{+12.86}{9.39}{16.37} & \pairedcell{+10.57}{7.08}{14.41} & \pairedcell{+7.86}{5.54}{10.28} \\[2pt]
\midrule
GPF & \pairedcell{+46.27}{42.36}{49.82} & \pairedcell{+48.25}{45.22}{51.53} & \pairedcell{+46.17}{43.71}{49.32} & \pairedcell{+37.48}{34.93}{39.92} & \pairedcell{+44.38}{40.82}{47.90} & \pairedcell{+47.77}{44.95}{50.72} & \pairedcell{+31.76}{28.23}{35.27} & \pairedcell{+21.33}{17.31}{25.89} & \pairedcell{+17.97}{15.46}{20.67} \\[2pt]
GPF-plus & \pairedcell{+26.26}{22.23}{30.20} & \pairedcell{+30.07}{28.04}{31.97} & \pairedcell{+20.22}{18.02}{22.39} & \pairedcell{+31.43}{28.61}{34.12} & \pairedcell{+32.03}{30.10}{33.81} & \pairedcell{+33.57}{32.22}{34.87} & \pairedcell{+22.48}{18.10}{26.98} & \pairedcell{+13.74}{12.07}{15.56} & \pairedcell{+11.61}{9.80}{13.53} \\[2pt]
\gpromptname & \pairedcell{+32.32}{28.93}{35.68} & \pairedcell{+42.41}{40.61}{44.20} & \pairedcell{+41.22}{39.64}{42.75} & \pairedcell{+26.26}{23.90}{28.62} & \pairedcell{+34.48}{31.80}{36.73} & \pairedcell{+36.31}{34.99}{37.65} & \pairedcell{+12.02}{7.33}{16.52} & \pairedcell{+8.27}{6.21}{10.36} & \pairedcell{+5.60}{3.49}{7.88} \\[2pt]
\midrule
\sharedname & \pairedcell{+31.70}{28.29}{35.00} & \pairedcell{+39.13}{36.60}{41.48} & \pairedcell{+37.15}{35.09}{39.23} & \pairedcell{+24.69}{22.70}{26.65} & \pairedcell{+35.52}{33.31}{37.55} & \pairedcell{+38.48}{37.22}{39.72} & \pairedcell{+5.44}{1.84}{9.07} & \pairedcell{+7.27}{4.64}{10.29} & \pairedcell{+4.10}{2.56}{5.58} \\[2pt]
\bottomrule
\end{tabular*}
\par\vspace{3pt}
\begin{minipage}{\linewidth}\fontsize{8}{9.5}\selectfont\raggedright \lang{Each cell: PRPrompt minus the reference (top); unadjusted 95\% paired support-bootstrap interval (below). All values preserve the original condition-wise statistics.}{每格上行为 PRPrompt 减去对照的差值，下行为未经多重比较校正的 95\% 配对支持 bootstrap 区间。数值沿用原始逐条件统计。}\end{minipage}
\end{minipage}\par

}]

\clearpage
\twocolumn[{
\section{Additional Datasets}
\label{app:extended}
Photo, Computers, Actor, Chameleon, and ogbn-arxiv use the same task-level statistics. All four methods are shown without averaging datasets. Original ProG sampling is retained. On ogbn-arxiv, seven nominal 3-shot and thirteen nominal 5-shot tasks lack the requested number of supports in a minority class, with actual $k_c$ between 1 and 4. Four shortage tasks have minimum count one and retain the inclusive prototype loss; otherwise LOO uses actual class counts. Pools are not replenished.
\par\medskip\noindent\begin{minipage}{\textwidth}
\centering
\linespread{1}\selectfont
\captionsetup{font=small,skip=5pt}
\captionof{table}{\lang{\textbf{Accuracy on additional datasets.}}{\textbf{扩展数据集的Accuracy。}}}
\label{tab:extended-accuracy}
\fontsize{9}{11}\selectfont
\setlength{\tabcolsep}{1pt}
\renewcommand{\arraystretch}{1.18}
\setlength{\heavyrulewidth}{0.8pt}
\setlength{\lightrulewidth}{0.4pt}
\begin{tabular*}{\linewidth}{@{\extracolsep{\fill}}ll*{4}{r}@{}}
\toprule
\lang{Dataset}{数据集} & Shot & \frozenname & \ftname & \sharedname & \textbf{\method} \\
\midrule
Photo & 1 & \meanstd{47.73}{5.38} & \secondmeanstd{50.87}{6.88} & \meanstd{34.10}{4.92} & \bestmeanstd{60.96}{8.88} \\
 & 3 & \meanstd{58.82}{4.41} & \bestmeanstd{74.78}{5.52} & \meanstd{49.55}{7.34} & \secondmeanstd{72.96}{4.71} \\
 & 5 & \meanstd{63.09}{4.16} & \bestmeanstd{80.68}{3.32} & \meanstd{52.46}{6.11} & \secondmeanstd{73.10}{5.28} \\
\midrule
Computers & 1 & \meanstd{38.13}{6.59} & \secondmeanstd{39.95}{9.06} & \meanstd{33.65}{6.68} & \bestmeanstd{45.60}{8.85} \\
 & 3 & \meanstd{47.56}{3.76} & \bestmeanstd{62.19}{5.57} & \meanstd{46.29}{6.22} & \secondmeanstd{57.53}{7.03} \\
 & 5 & \meanstd{52.41}{2.89} & \bestmeanstd{68.79}{4.47} & \meanstd{51.29}{3.53} & \secondmeanstd{63.12}{3.12} \\
\midrule
Actor & 1 & \secondmeanstd{20.61}{1.22} & \meanstd{20.28}{3.19} & \meanstd{20.20}{1.59} & \bestmeanstd{20.64}{2.30} \\
 & 3 & \secondmeanstd{21.14}{1.46} & \meanstd{20.87}{1.86} & \meanstd{20.42}{1.27} & \bestmeanstd{21.24}{1.40} \\
 & 5 & \secondmeanstd{21.12}{1.63} & \meanstd{20.97}{1.80} & \meanstd{20.98}{1.21} & \bestmeanstd{21.34}{1.69} \\
\midrule
Chameleon & 1 & \meanstd{21.47}{3.01} & \meanstd{20.91}{2.12} & \secondmeanstd{22.22}{3.11} & \bestmeanstd{22.90}{3.53} \\
 & 3 & \meanstd{22.85}{2.73} & \secondmeanstd{24.00}{3.44} & \meanstd{23.42}{2.76} & \bestmeanstd{25.09}{4.10} \\
 & 5 & \meanstd{23.71}{3.05} & \secondmeanstd{26.05}{3.35} & \meanstd{24.03}{3.34} & \bestmeanstd{27.29}{3.80} \\
\midrule
ogbn-arxiv & 1 & \meanstd{10.93}{1.66} & \bestmeanstd{16.28}{3.79} & \meanstd{10.57}{1.80} & \secondmeanstd{11.07}{1.67} \\
 & 3 & \meanstd{16.13}{2.03} & \bestmeanstd{21.67}{3.28} & \meanstd{18.06}{2.22} & \secondmeanstd{19.97}{2.26} \\
 & 5 & \meanstd{19.54}{1.99} & \bestmeanstd{24.03}{2.57} & \meanstd{22.52}{2.23} & \secondmeanstd{23.95}{1.94} \\
\bottomrule
\end{tabular*}
\par\vspace{3pt}
\begin{minipage}{\linewidth}\fontsize{8}{9.5}\selectfont\raggedright \lang{Local original-ProG runs with support/test overlap: mean $\pm$ sample std. over 20 tasks after averaging three training seeds per task. \textbf{Bold}/\underline{underline}: best/second in each condition. No cross-condition average.}{本地原ProG实现，保留支持／测试重叠：先平均每任务3个训练seed，再报告20个任务的均值 $\pm$ 样本标准差。\textbf{粗体}／\underline{下划线}表示同条件最优／次优；不跨条件平均。}\lang{ For ogbn-arxiv, shots are nominal: 7/20 three-shot and 13/20 five-shot tasks have a class shortage; see text.}{ogbn-arxiv为名义shot：3-shot中7/20、5-shot中13/20个任务存在少数类支持不足，见正文说明。}\end{minipage}
\end{minipage}\par

\par\medskip\noindent\begin{minipage}{\textwidth}
\centering
\linespread{1}\selectfont
\captionsetup{font=small,skip=5pt}
\captionof{table}{\lang{\textbf{Macro-F1 on additional datasets.}}{\textbf{扩展数据集的Macro-F1。}}}
\label{tab:extended-macro_f1}
\fontsize{9}{11}\selectfont
\setlength{\tabcolsep}{1pt}
\renewcommand{\arraystretch}{1.18}
\setlength{\heavyrulewidth}{0.8pt}
\setlength{\lightrulewidth}{0.4pt}
\begin{tabular*}{\linewidth}{@{\extracolsep{\fill}}ll*{4}{r}@{}}
\toprule
\lang{Dataset}{数据集} & Shot & \frozenname & \ftname & \sharedname & \textbf{\method} \\
\midrule
Photo & 1 & \meanstd{48.13}{4.86} & \secondmeanstd{50.48}{6.11} & \meanstd{33.00}{4.70} & \bestmeanstd{60.54}{7.59} \\
 & 3 & \meanstd{59.13}{3.63} & \bestmeanstd{73.90}{4.73} & \meanstd{49.41}{6.56} & \secondmeanstd{71.58}{3.59} \\
 & 5 & \meanstd{62.98}{3.93} & \bestmeanstd{79.62}{2.74} & \meanstd{51.20}{5.62} & \secondmeanstd{72.32}{4.45} \\
\midrule
Computers & 1 & \meanstd{36.39}{5.76} & \secondmeanstd{38.66}{8.09} & \meanstd{31.90}{5.77} & \bestmeanstd{44.77}{7.13} \\
 & 3 & \meanstd{44.93}{3.70} & \bestmeanstd{61.51}{5.38} & \meanstd{43.38}{5.89} & \secondmeanstd{54.79}{5.46} \\
 & 5 & \meanstd{48.50}{3.17} & \bestmeanstd{68.65}{3.88} & \meanstd{46.89}{4.24} & \secondmeanstd{59.44}{2.98} \\
\midrule
Actor & 1 & \secondmeanstd{19.12}{0.98} & \meanstd{16.16}{2.35} & \bestmeanstd{19.21}{0.98} & \meanstd{19.00}{1.64} \\
 & 3 & \meanstd{19.23}{1.62} & \meanstd{18.45}{1.20} & \secondmeanstd{19.65}{0.90} & \bestmeanstd{20.06}{0.97} \\
 & 5 & \meanstd{19.44}{1.57} & \meanstd{19.10}{0.98} & \secondmeanstd{20.16}{0.93} & \bestmeanstd{20.49}{1.51} \\
\midrule
Chameleon & 1 & \secondmeanstd{20.38}{2.89} & \meanstd{14.27}{2.60} & \meanstd{19.87}{3.30} & \bestmeanstd{20.72}{3.94} \\
 & 3 & \secondmeanstd{22.04}{2.64} & \meanstd{19.98}{3.68} & \meanstd{21.55}{3.10} & \bestmeanstd{23.32}{3.93} \\
 & 5 & \meanstd{22.75}{3.02} & \secondmeanstd{23.10}{3.33} & \meanstd{22.00}{3.37} & \bestmeanstd{26.03}{3.39} \\
\midrule
ogbn-arxiv & 1 & \secondmeanstd{8.35}{1.08} & \bestmeanstd{9.95}{1.45} & \meanstd{7.82}{0.98} & \meanstd{7.92}{1.05} \\
 & 3 & \meanstd{12.77}{1.08} & \bestmeanstd{14.46}{1.25} & \meanstd{13.08}{1.18} & \secondmeanstd{14.16}{1.05} \\
 & 5 & \meanstd{15.59}{1.01} & \secondmeanstd{16.68}{1.17} & \meanstd{16.44}{1.08} & \bestmeanstd{17.21}{0.99} \\
\bottomrule
\end{tabular*}
\par\vspace{3pt}
\begin{minipage}{\linewidth}\fontsize{8}{9.5}\selectfont\raggedright \lang{Local original-ProG runs with support/test overlap: mean $\pm$ sample std. over 20 tasks after averaging three training seeds per task. \textbf{Bold}/\underline{underline}: best/second in each condition. No cross-condition average.}{本地原ProG实现，保留支持／测试重叠：先平均每任务3个训练seed，再报告20个任务的均值 $\pm$ 样本标准差。\textbf{粗体}／\underline{下划线}表示同条件最优／次优；不跨条件平均。}\lang{ For ogbn-arxiv, shots are nominal: 7/20 three-shot and 13/20 five-shot tasks have a class shortage; see text.}{ogbn-arxiv为名义shot：3-shot中7/20、5-shot中13/20个任务存在少数类支持不足，见正文说明。}\end{minipage}
\end{minipage}\par

}]
\noindent Full fine-tuning outperforms \method{} on Photo and Computers at 3/5 shots, and the ogbn-arxiv 1-shot gap is particularly pronounced. Several Actor routing-effect intervals include zero; its 1-shot Macro-F1 mean difference from the shared prompt is negative. The 1-shot routing effects on Chameleon and ogbn-arxiv are also uncertain. These results do not establish universal superiority across graphs. All five metrics, paired intervals, and actual class-support counts remain in the data package.
